\begin{abstract}
Domain-generalizable person re-identification (DG-ReID) typically optimizes symmetric distances on Euclidean embeddings, yet probe-to-gallery retrieval in heterogeneous camera networks is not symmetric: viewpoint, resolution, and sensing pipelines induce direction-dependent mismatch that a single metric may discard. We study asymmetric, Randers/Finsler-style retrieval scores on top of alignment–uniformity training \cite{cho2024generalizable}, using a structured identity/drift representation while keeping the backbone geometry Euclidean. We show that nuisance-aligned drift objectives conflict with domain-invariant training pressures, and formulate a unified asymmetric scoring layer with norm-constrained drift that recovers Euclidean geometry at initialization. On M+MS+CS$\to$CUHK03, five-seed mean mAP is $43.8\pm0.3\%$; a protocol-matched run achieves $44.3\%$ mAP ($+1.9\%$ vs.\ a configuration-matched Euclidean baseline). A controlled corruption experiment isolates the mechanism: a pairwise monotonicity loss encodes corruption severity in a learned scalar (Spearman $\rho\!\approx\!{-0.91}$). The Randers correction produces the predicted directional distance gap and balanced evaluation at matched counts confirms that clean-probe retrieval is consistently easier than corrupted-probe retrieval. The results suggest that drift-subspace variance regularization, not explicit domain-mining auxiliaries, is the critical ingredient for stable asymmetric retrieval. Code is available at \url{https://github.com/leezhuole/Manifold-Aware-Person-Re-Identification}.
\end{abstract}